% Concept proposal, not yet inserted into the manuscript.
% Requires graphicx. Include from a document whose working directory is paper/.
\begin{figure*}[t]
  \centering
  \includegraphics[width=\textwidth]{figures/central-concept-v1/central-concept.pdf}
  \caption{\textbf{From a learned answer to reusable knowledge.}
  (a) A representative pre-norm Transformer layer combines contextual retrieval
  with a persistent MLP mapping, retaining both residual paths.
  (b) The one-hop query $q_2=(e_2,r_2)$ and the two-hop query
  $q_{12}=(e_1,r_1,r_2)$ request the same answer but can reach different
  retrieval states. Reusing an association requires access to the learned
  parameters and a compatible input; subsequent computation determines the
  behavioral effect. The illustrated activation strips are conceptual, not
  measured features or an embedding projection.
  (c) Editing the first fact changes the intermediate entity; editing the
  second tests reuse of an updated association across queries. The arrows
  indicate desired dependency changes, not guaranteed update propagation.
  (d) Entity-centered addressing and re-entering compatible computation are
  two routes by which knowledge can become reusable. Architectural access and
  handoff form are distinct requirements. The displayed update is one SGD
  step on the MLP down-projection; its direct readout identity fixes the
  features and does not describe the full effect of arbitrary multi-step
  training. The schematic does not assign knowledge and reasoning to fixed
  layer regions.}
  \label{fig:knowledge-reuse-concept-proposal}
\end{figure*}
